\documentclass[10pt,a4paper]{amsart}
\usepackage[margin=19mm]{geometry}
\usepackage{amsmath,amssymb,mathtools}
\usepackage[scr=rsfs,cal=euler]{mathalfa}
\usepackage{xfrac}
\usepackage[unicode,hidelinks,bookmarks=false]{hyperref}
\newcommand{\ie}{i.e.\ }
\newcommand{\cf}{cf.\ }
\newcommand{\resp}{respectively\ }
\newcommand{\Subsection}{\subsection}
\providecommand{\nobreakdash}{\nobreak\hbox{-}}
\setlength{\emergencystretch}{2em}
\multlinegap0pt
\usepackage{amssymb}
\usepackage{tikz}
\usepackage{mathtools}
\usepackage{csquotes}
\newtheorem{lemma}{Lemma}
\title{Maximising the number of solutions to linear equations}
\author{\v{Z}arko Ran\dj elovi\'c, Xuancheng Shao, Max Wenqiang Xu, Shengtong Zhang, Yuan Zhou}
\date{Source: arXiv:2609.06975v1 (CC BY 4.0)}
\begin{document}
\maketitle

\noindent\emph{Source and licence.} Maximising the number of solutions to linear equations. Version arXiv:2609.06975v1 (CC BY 4.0), distributed by arXiv under the Creative Commons Attribution 4.0 International licence, http://creativecommons.org/licenses/by/4.0/, as recorded in the arXiv licence metadata for this exact version and shown on the abstract page https://arxiv.org/abs/2609.06975v1. Full licence notice: \texttt{licenses/arxiv-2609.06975-CC-BY-4.0.txt}.\par\smallskip

\noindent\textit{Editorial scope.} Counted unit: the article's lemma ABAcupB (source0.tex lines 289--293). The article's complete original proof of it is delivered with it (source0.tex lines 1183--1477, one \texttt{proof} environment of 295 lines, which the article places away from the statement). No mathematics is written, repaired or re-derived: the statement and the proof are the article's own lines, in the article's own notation, and no retained line is substituted or re-flowed.  Retained uncounted as the source-local context the proof needs: lines 272--276.  Nothing was omitted from any retained span, so the package carries no omission record.  No retained line contains a citation, so the wrapper carries no bibliography and no bibliography entry.  No retained line contains a figure environment, an image-inclusion command or drawing code, so no image is delivered and the package declares no asset dependency.  Editorial formatting only, in addition: an AMS \texttt{amsart} preamble that restates the article's own declarations and macros used by the retained lines, namely the environments \texttt{lemma} on separate counters, together with the standard packages those lines need.  The retained environments print in this document as Lemma 1 in order of appearance, whereas the article prints the same items with the numbering of its own full document.  The complete native source of the article as distributed in the arXiv e-print for this exact version is bundled unchanged as \texttt{sources/209595/source0.tex}.
\begin{lemma}
    If $|A| + |C| \geq |B|$ and $|B| + |C| \geq |A|$, then
    \label{lemma ABC}
    $$T(A, B, C) \leq |A||B| - \frac{1}{4} \max(|A| + |B| - |C|, 0)^2 + 1.$$
\end{lemma}

\begin{lemma}
\label{lemma ABAcupB}
For any set $A$ consisting of integers congruent to $1$ modulo $3$, and $B$ consisting of integers congruent to $2$ modulo $3$, we have 
$$T(A, B, 3 \cdot (A \cup B)) \leq \frac{6}{25} \left(|A| + |B|\right)^2 + 2.$$
\end{lemma}

\begin{proof}[Proof of Lemma \ref{lemma ABAcupB}]
Write
\[
a:=|A|,\qquad b:=|B|.
\]
Since every element of $A$ is congruent to $1 \pmod 3$ and every element of $B$ is congruent to $2 \pmod 3$, every element of $3(A\cup B)$ is congruent to $3$ or $6 \pmod 9$. In particular,
\begin{equation}
\label{eq:no9}
3(A\cup B)\cap 9\mathbb Z=\varnothing.
\end{equation}

For $r\in\{1,4,7\}$ and $s\in\{2,5,8\}$, define
\[
A^{(r)}:=\{x\in A:x\equiv r \pmod 9\},\qquad
B^{(s)}:=\{y\in B:y\equiv s \pmod 9\}.
\]
Thus $\{A^{(1)},A^{(4)},A^{(7)}\}$ partitions $A$ and $\{B^{(2)},B^{(5)},B^{(8)}\}$ partitions $B$.

We say that a residue class is \emph{popular} in $A$ if it contains at least $\frac45a$ elements of $A$, and similarly for $B$.

\medskip
\noindent\textbf{Case 1: neither $A$ nor $B$ has a popular residue class.}

Then
\[
|A^{(r)}|<\frac45a\quad (r\in\{1,4,7\}),\qquad
|B^{(s)}|<\frac45b\quad (s\in\{2,5,8\}).
\]
Let $D$ denote the number of pairs $(x,y)\in A\times B$ such that $x+y\equiv 0\pmod 9$. It is clear that
\[
D=|A^{(1)}||B^{(8)}|+|A^{(4)}||B^{(5)}|+|A^{(7)}||B^{(2)}|.
\]
Every such pair satisfies $x+y\in 9\mathbb Z$, so by \eqref{eq:no9} it contributes nothing to $T(A,B,3(A\cup B))$. Hence
\begin{equation}
\label{eq:TleabminusD}
T(A,B,3(A\cup B))\le ab-D.
\end{equation}

We claim that
\[
D\ge \frac1{25}ab.
\]
Indeed, set
\[
\alpha_1:=\frac{|A^{(1)}|}{a},\quad
\alpha_2:=\frac{|A^{(4)}|}{a},\quad
\alpha_3:=\frac{|A^{(7)}|}{a},
\]
and
\[
\beta_1:=\frac{|B^{(8)}|}{b},\quad
\beta_2:=\frac{|B^{(5)}|}{b},\quad
\beta_3:=\frac{|B^{(2)}|}{b}.
\]
Then $\alpha_i,\beta_i\ge 0$, $\alpha_1+\alpha_2+\alpha_3=1$, $\beta_1+\beta_2+\beta_3=1$, and $\alpha_i,\beta_i\le \frac45$ for all $i$. Also
\[
\frac{D}{ab}=\alpha_1\beta_1+\alpha_2\beta_2+\alpha_3\beta_3.
\]
After relabelling the indices, we may assume
\[
\beta_1\le \beta_2\le \beta_3.
\]
For fixed $(\beta_1,\beta_2,\beta_3)$, the minimum of $\alpha_1\beta_1+\alpha_2\beta_2+\alpha_3\beta_3$ subject to
\[
\alpha_i\ge 0,\qquad \alpha_1+\alpha_2+\alpha_3=1,\qquad \alpha_i\le \frac45
\]
is attained at
\[
\alpha_1=\frac45,\qquad \alpha_2=\frac15,\qquad \alpha_3=0.
\]
Therefore
\[
\frac{D}{ab}\ge \frac45\beta_1+\frac15\beta_2\ge \frac15(\beta_1+\beta_2).
\]
Since $\beta_3\le \frac45$, we have $\beta_1+\beta_2=1-\beta_3\ge \frac15$, and so
\[
\frac{D}{ab}\ge \frac1{25}.
\]
This proves the claim.

Combining the claim with \eqref{eq:TleabminusD}, we obtain
\[
T(A,B,3(A\cup B))\le ab-\frac1{25}ab=\frac{24}{25}ab
\le \frac{24}{25}\cdot \frac{(a+b)^2}{4}
=\frac6{25}(a+b)^2.
\]

\medskip
\noindent\textbf{Case 2: at least one of $A$ and $B$ has a popular residue class.}

By multiplying all elements of $A\cup B$ with $2$ if needed, we may assume that $A$ has a popular residue class. Choose $r_0\in\{1,4,7\}$ such that
\[
|A^{(r_0)}|\ge \frac45a.
\]
Set
\[
A^{\star}:=A^{(r_0)},\qquad A^{\circ}:=A\setminus A^{\star},
\qquad a_{\star}:=|A^{\star}|,\qquad a_{\circ}:=|A^{\circ}|.
\]
Then
\begin{equation}
\label{eq:astar}
a_{\star}\ge \frac45a,\qquad a_{\circ}\le \frac15a.
\end{equation}

Let $s_0\in\{2,5,8\}$ be the unique residue class such that
\[
r_0+s_0\equiv 3\pmod 9.
\]
Set
\[
B^{\star}:=B^{(s_0)},\qquad B^{\circ}:=B\setminus B^{\star},
\qquad b_{\star}:=|B^{\star}|,\qquad b_{\circ}:=|B^{\circ}|.
\]

If $x\in A^{\star}$ and $y\in B^{\star}$, then $x+y\equiv 3\pmod 9$, hence any solution
\[
x+y=z\in 3(A\cup B)
\]
must satisfy $z\in 3A$. Likewise, if $x\in A^{\star}$ and $y\in B^{\circ}$, then $x+y\not\equiv 3\pmod 9$, so any solution with $z\in 3(A\cup B)$ must satisfy $z\in 3B$ (if $x+y\equiv 0\pmod 9$, then there is no solution by \eqref{eq:no9}). Therefore
\begin{equation}
\label{eq:decomp}
T(A,B,3(A\cup B))
=
T(A^{\star},B^{\star},3A)
+
T(A^{\star},B^{\circ},3B)
+
T(A^{\circ},B,3(A\cup B)).
\end{equation}
Also
\begin{equation}
\label{eq:lasttrivial}
T(A^{\circ},B,3(A\cup B))\le a_{\circ}b.
\end{equation}

We now distinguish subcases.

\medskip
\noindent\textbf{Subcase 2a: $b_{\star}>a_{\star}+a$.}

Then
\[
b\ge b_{\star}>a_{\star}+a\ge \frac45a+a=\frac95a.
\]
Hence
\[
T(A,B,3(A\cup B))\le ab\le \frac6{25}(a+b)^2,
\]
as $b\ge \frac95a\ge\frac32a$.

\medskip
\noindent\textbf{Subcase 2b: $b_{\star}\le a_{\star}+a$ and $a_{\star}\le b_{\circ}+b$.}

Then Lemma~\ref{lemma ABC} applies to both $T(A^{\star},B^{\star},3A)$ and $T(A^{\star},B^{\circ},3B)$. Thus
\[
T(A^{\star},B^{\star},3A)
\le
a_{\star}b_{\star}
-\frac14\max(a_{\star}+b_{\star}-a,0)^2
+1,
\]
and
\[
T(A^{\star},B^{\circ},3B)
\le
a_{\star}b_{\circ}
-\frac14\max(a_{\star}+b_{\circ}-b,0)^2
+1.
\]
Combining these with \eqref{eq:decomp} and \eqref{eq:lasttrivial},
\[
T(A,B,3(A\cup B))
\le
ab
-\frac14(u^2+v^2)
+2,
\]
where
\[
u:=\max(a_{\star}+b_{\star}-a,0),\qquad
v:=\max(a_{\star}+b_{\circ}-b,0).
\]
Since $b_{\star}+b_{\circ}=b$, we have
\[
(a_{\star}+b_{\star}-a)+(a_{\star}+b_{\circ}-b)=2a_{\star}-a.
\]
Hence
\[
u+v
\ge
\max\bigl((a_{\star}+b_{\star}-a)+(a_{\star}+b_{\circ}-b),\,0\bigr)
=
2a_{\star}-a.
\]
Using \eqref{eq:astar}, this gives
\[
u+v\ge 2a_{\star}-a\ge \frac35a.
\]
Therefore
\[
u^2+v^2\ge \frac12(u+v)^2\ge \frac12\left(\frac35a\right)^2=\frac{9}{50}a^2,
\]
and so
\[
T(A,B,3(A\cup B))
\le
ab-\frac14\cdot \frac{9}{50}a^2+2
=
ab-\frac{9}{200}a^2+2.
\]
Finally, it is easy to check that 
\[
ab-\frac{9}{200}a^2\le \frac6{25}(a+b)^2
.\]
Thus
\[
T(A,B,3(A\cup B))\le \frac6{25}(a+b)^2+2.
\]

\medskip
\noindent\textbf{Subcase 2c: $b_{\star}\le a_{\star}+a$ and $a_{\star}>b_{\circ}+b$.}

In this subcase Lemma~\ref{lemma ABC} still applies to $T(A^{\star},B^{\star},3A)$, so
\[
T(A^{\star},B^{\star},3A)
\le
a_{\star}b_{\star}
-\frac14\max(a_{\star}+b_{\star}-a,0)^2
+1.
\]
For the second term we use the trivial bound
\[
T(A^{\star},B^{\circ},3B)\le b_{\circ}b.
\]
Together with \eqref{eq:decomp} and \eqref{eq:lasttrivial}, this gives
\[
T(A,B,3(A\cup B))
\le
a_{\circ}b
+
a_{\star}b_{\star}
-\frac14\max(a_{\star}+b_{\star}-a,0)^2
+1
+
b_{\circ}b.
\]
Since $b_{\star}=b-b_{\circ}$ and $a_{\star} + b_{\star} - a = b_{\star} - a_{\circ}$, we obtain
\[
T(A,B,3(A\cup B))
\le
ab
-
(a_{\star}-b)b_{\circ}
-
\frac14\max(b_{\star}-a_{\circ},0)^2
+1.
\]
Discarding the nonnegative term $(a_{\star}-b)b_{\circ}$, we get
\begin{equation}
\label{eq:subcase2c-main}
T(A,B,3(A\cup B))
\le
ab
-
\frac14\max(b_{\star}-a_{\circ},0)^2
+1.
\end{equation}

Now $a_{\star}>b_{\circ}+b$ implies
\[
b_{\star}-a_{\circ}
=
b-b_{\circ}-a+a_{\star}
>
2b-a.
\]
It follows that
\[
\max(b_{\star}-a_{\circ},0)\ge \max(2b-a,0).
\]
So \eqref{eq:subcase2c-main} yields
\[
T(A,B,3(A\cup B))
\le
ab-\frac14\max(2b-a,0)^2+1.
\]
We claim that
\[
ab-\frac14\max(2b-a,0)^2\le \frac6{25}(a+b)^2.
\]
Indeed, setting $x = a / b$, this is equivalent to the elementary inequality
$$x - \frac{1}{4}\max\left(2 - x, 0\right)^{2} \leq \frac{6}{25}\left(x + 1\right)^{2}$$
which holds for any $x \in \mathbb{R}$. This completes the proof.
\end{proof}

\end{document}
